\documentclass[12pt]{article}
\usepackage{amsmath, amssymb}
\usepackage[margin=1in]{geometry}
\setlength{\parskip}{6pt}
\title{QUBO Formulation: Quantum Circuit Qubit Routing}
\date{}
\begin{document}
\maketitle

\subsection*{Quantum Circuit Qubit Routing}

\paragraph{Problem Statement.}
Given a quantum hardware coupling graph and an initial qubit placement, a set of $K$ candidate routing sequences is pre-generated. Each candidate sequence $k$ incurs a known swap cost $c_k$, representing the number of two-qubit swaps required, and a nonadjacency penalty $p_k$, which is zero if the target qubits become adjacent after applying the sequence and a large constant otherwise. The objective is to select exactly one candidate sequence such that the sum of its swap cost and nonadjacency penalty is minimized.

\paragraph{Decision Variables.}
For each candidate sequence $k \in \{1, \dots, K\}$, we introduce a binary decision variable $y_k \in \{0,1\}$. The variable $y_k$ equals $1$ if and only if candidate sequence $k$ is selected, and $0$ otherwise.

\paragraph{QUBO Derivation.}
\textbf{Step 1 -- Objective Function.} The original minimization objective is expressed as
\begin{align}
\min_{y} \sum_{k=1}^{K} (c_k + p_k) y_k.
\end{align}
Since QUBO solvers minimize energy, we retain this form directly.

\textbf{Step 2 -- Constraint.} The requirement to select exactly one sequence is written as
\begin{align}
\sum_{k=1}^{K} y_k = 1.
\end{align}

\textbf{Step 3 -- Penalty Term Expansion.} We enforce the constraint by adding a quadratic penalty term with weight $\lambda > 0$:
\begin{align}
P(y) = \lambda \left( \sum_{k=1}^{K} y_k - 1 \right)^2.
\end{align}
Expanding the square yields
\begin{align}
P(y) = \lambda \left( \sum_{k=1}^{K} y_k^2 - 2 \sum_{1 \le i < j \le K} y_i y_j + 1 \right).
\end{align}
Using the idempotent property of binary variables, $y_k^2 = y_k$, we simplify the expression to
\begin{align}
P(y) = \lambda \sum_{k=1}^{K} y_k - 2\lambda \sum_{1 \le i < j \le K} y_i y_j + \lambda.
\end{align}

\textbf{Step 4 -- Combined Hamiltonian.} Adding the objective and the penalty term gives the full QUBO Hamiltonian $H(y)$:
\begin{align}
H(y) = \sum_{k=1}^{K} (c_k + p_k) y_k + \lambda \sum_{k=1}^{K} y_k - 2\lambda \sum_{1 \le i < j \le K} y_i y_j + \lambda.
\end{align}
Grouping linear terms, we obtain
\begin{align}
H(y) = \sum_{k=1}^{K} (c_k + p_k + \lambda) y_k - 2\lambda \sum_{1 \le i < j \le K} y_i y_j + \lambda.
\end{align}

\textbf{Step 5 -- Q Matrix and Offset.} Reading off the coefficients from the standard upper-triangular QUBO representation $H(y) = \sum_{k=1}^{K} Q_{k,k} y_k + \sum_{1 \le i < j \le K} Q_{i,j} y_i y_j + c$, we obtain:
\begin{align}
Q_{k,k} &= c_k + p_k + \lambda, \quad \text{for } k \in \{1, \dots, K\}, \\
Q_{i,j} &= -\lambda, \quad \text{for } 1 \le i < j \le K, \\
c &= \lambda.
\end{align}
Note that the Hamiltonian contains the term $-2\lambda y_i y_j$ for each pair, which corresponds to $Q_{i,j} = -\lambda$ in the standard QUBO convention.

\paragraph{Penalty Weight Justification.}
The penalty weight $\lambda$ must be chosen sufficiently large to ensure that any constraint violation increases the total energy above the minimum achievable by a feasible solution. The maximum possible reduction in the objective term $\sum_{k=1}^{K} (c_k + p_k) y_k$ occurs when switching from an infeasible assignment to the best feasible assignment. Since at most one variable can change from $1$ to $0$ (or vice versa) when moving between feasible and infeasible regions, the maximum objective gain from violating the constraint is bounded by $\max_{k} (c_k + p_k)$. Therefore, setting
\begin{align}
\lambda > \max_{k \in \{1, \dots, K\}} (c_k + p_k)
\end{align}
guarantees that the penalty incurred by any infeasible configuration strictly exceeds any potential reduction in the swap cost or penalty terms.

\paragraph{Correctness Argument.}
Minimizing $H(y)$ over $\{0,1\}^K$ recovers the optimal routing sequence because (i) all feasible solutions satisfy $\sum_{k=1}^{K} y_k = 1$, causing the penalty term to evaluate to exactly zero, so their energy equals the original objective value plus the constant offset $\lambda$; and (ii) any infeasible solution violates the constraint, incurring a penalty of at least $\lambda$. By the choice of $\lambda$, this penalty strictly dominates any possible decrease in the linear objective terms, ensuring that the global minimum of $H(y)$ must occur at a feasible point. Among feasible points, $H(y)$ reduces to the original objective function, so the minimizer corresponds exactly to the candidate sequence with the minimum total cost.

\paragraph{Illustrative Example.}
Consider an instance with $K=3$ candidate sequences. Let the instance parameters be $c_1 = 2, p_1 = 0$; $c_2 = 3, p_2 = 0$; $c_3 = 4, p_3 = 50$. We choose $\lambda = 51$, which satisfies $\lambda > \max\{2, 3, 54\}$. Substituting these values into the general formulas yields the concrete QUBO:
\begin{align}
Q_{1,1} &= 2 + 0 + 51 = 53, \\
Q_{2,2} &= 3 + 0 + 51 = 54, \\
Q_{3,3} &= 4 + 50 + 51 = 55, \\
Q_{1,2} = Q_{1,3} = Q_{2,3} &= -51, \\
c &= 51.
\end{align}
The resulting Hamiltonian is
\begin{align}
H(y) = 53 y_1 + 54 y_2 + 55 y_3 - 51 y_1 y_2 - 51 y_1 y_3 - 51 y_2 y_3 + 51.
\end{align}
Evaluating $H(y)$ over the feasible space $\{(1,0,0), (0,1,0), (0,0,1)\}$ gives energies $104$, $105$, and $106$, respectively. Infeasible configurations, such as $(1,1,0)$, yield $H(1,1,0) = 53 + 54 - 51 + 51 = 107$, which is strictly higher. The global minimum occurs at $y^* = (1,0,0)$ with energy $104$, correctly selecting candidate sequence $1$ as the optimal routing sequence.
\end{document}